\section[后记\q 专访“拿铁因子”概念创始人]{后记\\ 专访“拿铁因子”概念创始人}

{\heiti 就我所知，柔伊与病床上母亲的对话，就是关于如何不要生有遗憾的那一段，是基于您和您的祖母萝丝·巴哈临终前的对话。请问，您愿意多分享一些祖母的故事，以及她如何影响了您和您的人生吗？} \vspace{1em}

我的祖母萝丝是一个不同凡响的女人。她30岁时，做了一个决定：那就是她不要再过贫穷的生活。这个决定改变了我们一家。

她当时住在威斯康辛州的密尔沃基，在盖伯尔百货公司卖假发。她和我的祖父杰克都没上过大学，他们是典型的美国中西部劳工阶级，薪水刚好支付支出，生活拮据。但是我祖母却有一个梦想，就是哪天能够过上富裕的生活。因此，她在30岁生日时，对祖父说：“是时候改变我们的生活了，我们来存钱吧！”

后来，她就和祖父一起，一周存1美元。我说真的，就是1美元。之后，她开始带午餐去上班。百货公司的同事都取笑她：“萝丝，你也太省了。和我们一起去用餐吧！”祖母说，一开始时，她也很难过。但是，她知道她省钱的原因。她想离开密尔沃基，她想去一个冬天不会太冷的地方，老了在温暖的地方退休。

后来，祖母用了一生的时间，把日用节省下来去投资，积蓄累积到百万美元。她也把她的理财秘笈传给了我。我7岁时，她帮我买了第一支股票，当时我们正在麦当劳，7岁的我觉得麦当劳是世上最好的餐厅。祖母成为我的第一个理财导师。而她的教导，改变了我的人生。

我的第一本著作《聪明女人致富七招》（\textit{Smart Women Finish Rich}）在20年前出版，就是把我从祖母那里学来的秘笈分享给读者。我是从1997年开始动手写的，祖母知道我的写作计划。可是，86岁的老祖母却中风了。当时，祖父已经过世10年，祖母寡居时，1天要走8公里，1杯果菜汁，1周出门约会3次，而且是不同的男性朋友，当然，我们是在丧礼上才知道的。

祖母中风之后，我们把她从加州的拉古纳海湾移去了湾区的一家护理中心，离我上班的地方和住家都非常近。这样，我就可以每天探望她。我还记得祖母最后的日子，我们的交互，就像昨天的事。

有一天，我问祖母，还有什么要教导我的。她想了一会儿，说道：“不，该告诉你的我都说了。你的人生会很棒的。”

之后，我又问她，有没有觉得缺憾的地方。她又想了想，说道：“没有，此生无憾。”然后，她就回忆了一生中值得感恩的人和事。

第二天一早，我又去看望祖母，问她：“晚上睡得好吗？”

“太不好了。”她说。“我整晚没睡，一直想今生有没有什么遗憾的地方。也谢谢你，提出了这个问题。”

我们一起大笑了。然后，她就拉着我的手，说起了她的憾事，有些事还是少女时代的事。然后她对我说：“大卫，现在我要你好好听着，我的遗憾不重要，重要的是，为什么会有这些遗憾。

“在我分享给你的每一个片刻，其实我就面临着人生的十字路口，我必须做一个选择：不然，我就选择一条安全的路，我知道结果会是什么；另一条是危险的路，我知道那条路的尽头有金子在等待我，但是我不知道走这条路要付出什么。”

“然后，我选择了另一条路吗？去冒险取得金子吗？但是我每次都选择了安全的道路。这是我后悔的事。现在我躺在这里，日子不多了，我不知道如果另做他选，会有什么结果。”

“但是，奶奶，”我说，“你的人生已经很棒了。”

尽管祖母体力衰微，但是她依然用力握住了我的手，说：“大卫，不要等到老了再来遗憾。冒险去吧！

“特别要记得，”奶奶补充说。“当你走到十字路口的时候，你会听见两种声音。一个成熟大男孩的声音‘大卫，选这条安全的道路。’而另一个小男孩会说：‘选这条。选这条。好有趣。我们来试试！’小男孩的声音比较小，但是很有激情，很想尝试，要给小男孩机会。并且分享给你的朋友。”

这是我最后一次去探望祖母。

我开车回公司，把车停在地下室，大哭起来。哭完，看着车子的后视镜，对自己说：“就到此为止了。我不能只做大公司的财务分析师过一辈子。我要把《聪明女人致富七招》写出来，帮助其他的人，一起过上富裕的生活。我虽然还不知道怎么着手，但是我对奶奶发誓，3年内我要离开这里，完成自己的梦想。我要让我的内心小男孩到舞台上来做主角。”

事实上，我用了4年，而不只是3年。但还是完成了。4年后，我离开了任职的大公司，从加州移居纽约，虽然我的亲戚和朋友都在加州。我去纽约，希望写更多理财书，帮助千万的读者。

这不是一件容易的事，应该说，困难重重。中间的辛苦，我就不说了。但我从来没有后悔。且果实甘甜。《聪明女人致富七招》跃身为畅销书，至今销售了百万册。出版5年后，我上了欧普拉的节目，面对观众说明我的理财观。

而我想要分享理财观念的期待，一直持续至今，写出了这本书。

说来，拿铁因子的说法是在欧普拉的节目上提出来的。我想写一本小书，即便是没习惯买财经书看的读者，也可以几个小时就翻完。可是我的出版社却不认可。后来，前前后后出版了7本书，出版社依然不看好拿铁因子。后来，我决定好歹要写出来，找新的出版社合作吧！所以，就有了现在这本书了。

柔伊的妈妈对她说的“要过尽兴的生活”这句话，成为柔伊的动力，让柔伊不再生活在小心翼翼中，不小看自己，要活出富足的人生。这里的富足，与金钱无关，而是精神层面的。

拿铁因子无关咖啡，也无关金钱。而是鼓励，大家勇于实现梦想。

别让人生留下遗憾。让你心中的小男孩，小女孩走到台前来。现在就开始吧！


\vspace{2em}

{\heiti 请问，您有没有写下一本书的打算？愿意透露给我们一些吗？}

\vspace{1em}

我过去一直想写一本书，并且对我和我的妻子说，一定会做到的，之后，也不会再写另一本了。虽然我的妻子完全不理会我说的。现在这本书完成了，尘埃落定了。其实通常这种时候，一个新点子就会蹦出来。

所以如果要简短回答，那就是会。我确实在准备新书，并且希望新书有趣，而且非常不同。

《拿铁因子》出版90天后，我将和家人去意大利佛罗伦斯，在那里度过一年。我确实过了没有遗憾的生活，而且希望我的孩子可以在意大利生活一段时间，可以在上大学开始他们的人生前，去各地看看。因此，如果一切顺利，在读者看到这本书时，我应该在意大利，吃着披萨，喝着葡萄酒，和妻子和孩子们在贡多拉上逍遥度日。像柔伊一样，休一个高调的休假研习。

可能我会在那里更新博客。大家如果愿意，欢迎来网站加入我们：\href{https://davidbach.com/}{\nolinkurl{www.davidbach.com}}。

\vspace{2em}

{\heiti 本书末你加入图表和数据，愿意就这个方面多聊聊吗？}\vspace{1em}

很高兴你看到了本书的亮点，我最喜欢的也是复利神奇力量的图表。“金钱的时间价值”在我20来岁时，就激励了我。我在第209页的表格上，显示出不同利息值下，金钱随着时间的增加。

后面一张表格曾经让20岁的我大彻大悟：那就是1926年以来的投资报酬率。不要相信，投资股票收不到回报之类的话。

再后面的空白表格是专为各位读者设计的，帮助落实你的理财计划。第一张表格题为拿铁因子的挑战。可以照本书建议的，纪录一天的开支，找出你自己的拿铁因子。\textbf{只需要一天的时间}。这一天不需要做任何改变，好像过去平常的每一天一样，平常怎么用钱的这一天就依然怎么用。之后，回想一天的开支，想出哪些是可以省下来的小钱，所谓的拿铁因子。计算一下金额会是多少，想一想，如果只要稍做改变，会有多大成效。

之后的一天，我们来加倍拿铁因子挑战。这个问题我们在书里没有讨论，但是其实很简单：写下一个月的开支，找出其中可以省下来的不重要开支。


\vspace{2em}

{\heiti 如果一位读者接触到您的第一本书就是《拿铁因子》，您会推荐的第二本书是哪一本？}\vspace{1em}

当然是《自动千万富翁》。这是我推荐大家都要看的书，也是目前最畅销的书。好读，又容易做到，不需要做预算，也可以成为百万富翁。另外，如果是女性读者，我会建议《聪明女人致富七招》，如果是夫妻一起读，我推荐《聪明夫妇一起致富》（\textit{Smart Couples Finish Rich}）。这些书可以帮助读者厘清自己的价值、人生的目标，而拟定一个适合于自己，切实可行的财务规画。这些部分在英文的新版本里，都做了增补。有兴趣的读者，可以找来参照。

\vspace{2em}

{\heiti 我当然希望我在20年前看过这本书。但是现在我已经四、五十岁了，拿铁因子还有用吗？}\vspace{1em}

答案是肯定的。只要开始，就会有收获。

就这方面，我在另一本书上做了讨论，这本书叫《起步晚，照样致富》（\textit{Start Late, Finish Rich}）。这本书就是为这样一群人写的，他们之前存的不多，或者贷款很多，可能还面临着一些意料之外的挑战。\textbf{起步晚，不等于未来就不能改变。梦想晚一点实现，也比不实现好}。

存钱和投资的要点是，金钱不知道他的主人几岁。这些只是数学。我们现在举一个例子，假设一位已经的50岁的人，有可能1天存在10美元吗？那么，另外一半可能也存10美元吗？那两个人的家庭1天就可以存20美元了，1年就是7,300美元了。如果用来投资，回报率用10％来算，20年后就是461,696美元，差不多是50万美元了。

如果可以1天存40美元，1年就是14,600美元，以10％报酬率来算，20年就有将近100万美元。

最重要的就是选一个可行的方案，就这样，不用担心其他。千里之行，始于足下。

\vspace{2em}

{\heiti 还有什么想对读者说的吗？如果读者想和您联系，可以如何找到您的联系信息？}\vspace{1em}

我最想对读者说的一句话是，感谢。感谢您花时间在这本书上。

我想从柔伊身上看到的是，这趟冒险之旅不是一个人的独行侠之旅。柔伊从她的咖啡师，从她的老板，从她的新朋友身上找到了鼓励和支持的力量。我鼓励大家也找到这些支持的力量。

有个很好的办法是大家可以创建一些拿铁俱乐部，一起来讨论，一起来督促实践。比方说，一起讨论梦想，一起讨论如何实现它。如果可以和朋友一起实现更富足生活的模式，生活就会有所改变。

也欢迎大家来逛一逛我的网页：\href{https://davidbach.com/}{\nolinkurl{www.davidbach.com}}。这里有更多的拿铁因子新观念，有许多有趣的新点子。我突然想到了什么，就会即时更新在网络上，而且是免费的，买书的钱也可以成为拿铁因子了。当然，如果集结成册，再出书的时候，还希望大家支持。

希望这本小书可以带来改变。也祝福各位读者，都可以实现自己的梦想。
